\chapter{Whittaker}\label{ch:3}
% 原PDF第19页
本章把视角从 $\Bun_n$ 上的特征层转向 Whittaker 模型。函数论陈述是：给定球面 Hecke 特征值后，归一化的 Whittaker 向量唯一。其几何体现是 Satake 范畴的 Whittaker 模型；对 $\GL_n$，它恢复了由局部系统构造的尖点自守层\cite{FGKV98,Gai18}。
\section{回顾：Hecke 特征层与尖点性}\label{sec:3.1}\index{尖点性|hyperpage}\index{常数项函子|hyperpage}
令 $\Bun_n:=\Bun_{\GL_n}$。对 $0\le d\le n$，第 $d$ 个微小 Hecke 对应为
\origdiagram{p019-hecke.pdf}{0.3871}
一个 $S$-点是元组
\[(x,M,M',M'\subset M\subset M'(x)),\]
其中 $M$ 与 $M'$ 是秩 $n$ 向量丛，$M/M'$ 在 $x$ 的图像上有限局部自由，秩为 $d$。归一化的 Hecke 函子为
\[H_n^d(\cK):=(\supp,\overrightarrow h)_!\overleftarrow h^*\cK\left(\frac{d(n-d)}2\right)[d(n-d)]\in D(X\times\Bun_n).\]
对 $X$ 上的局部系统 $E$，若
\[H_n^d(\cK)\simeq\bigwedge^dE\boxtimes\cK,\]
则称 $\cK$ 满足以 $E$ 为特征值的第 $d$ 个 Hecke 特征性质。
如\chref{2}所述，若 $\cK$ 是反常层，并满足第一特征性质以及迭代修改通常的对称相容性，则它满足所有微小特征性质。

这里使用的主要存在性定理是 Frenkel–Gaitsgory–Vilonen 定理\cite{FGKV98}：若 $E$ 几何不可约，则 $\Bun_n$ 上存在尖点 Hecke 特征层 $\Aut_E$，且它在每个连通分支 $\Bun_n^d$ 上的限制不可约。

还要回顾尖点性的几何定义。设 $G$ 约化，$P\subsetneq G$ 是真抛物子群，其幂幺根基为 $U$，Levi 商为 $M=P/U$。则
\origdiagram{p019-constant.pdf}{0.3575}
定义常数项函子
\[\CT_!:=q_!p^*.\]
若对每个真抛物子群 $P$ 都有 $\CT_!(\cK)=0$，则称层 $\cK\in D(\Bun_G)$ 为尖点层。
\section{经典 Hecke 特征函数}\label{sec:3.2}
\srcpage{20}
以下函数论段落假设 $k$ 有限。令 $F=k(X)$，$\A$ 为 $F$ 的阿代尔环，并置 $G=\GL_n$。设
\[\sigma:\Gal(\overline F/F)\longrightarrow\GL_n(\mathbb Q_\ell)\]
处处非分歧，等价地，它是 $X$ 上秩 $n$ 的局部系统。球面自守函数是定义在
\[\GL_n(F)\backslash\GL_n(\A)/\GL_n(\cO)\]
上的函数。对闭点 $x\in|X|$，记 $K_x$ 为局部域，$\cO_x$ 为其整数环，$q_x=\#k(x)$，并选取素元 $t_x$。置
\[M_n^d(\cO_x):=\GL_n(\cO_x)t_x^{(1^d,0^{n-d})}\GL_n(\cO_x)\subset\GL_n(K_x)\subset\GL_n(\A).\]
未归一化的局部 Hecke 算子为
\[T_x^df(g)=\int_{M_n^d(\cO_x)}f(gh)\,dh.\]
若对每个 $x\in|X|$ 都有
\[T_x^df=q_x^{d(n-d)/2}\Tr(\bigwedge^d\sigma(\Fr_x))f,\]
则称 $f$ 满足关于 $\sigma$ 的第 $d$ 个 Hecke 特征性质。
等价地，令 $\gamma_x$ 为由
\[\Tr(\bigwedge^d\gamma_x)=q_x^{d(n-1)/2}\Tr(\bigwedge^d\sigma(\Fr_x)),\qquad0\le d\le n\]
刻画的归一化 Satake 参数。采用 Tate 扭转的约定时，记为 $\gamma_x=\sigma((n-1)/2)(\Fr_x)$。于是，同一特征值为
\[q_x^{-d(d-1)/2}\Tr(\bigwedge^d\gamma_x).\]

\begin{transnote}{含糊点拨}
$\gamma_x$ 与 $\sigma(\Fr_x)$ 只差一个标量：$(n-1)/2$ 是半和正根 $\rho$ 在基本余权 $\varpi_1$ 上取的值，也就是 Borel 子群模特征 $\delta_B$ 的平方根在 $t^{\varpi_1}$ 处的值 $q_x^{(n-1)/2}$。这正是 \S\ref{sec:3.3} 用 $\rho$-移位（「上点」$\dot G$）所做的事：把 Whittaker 特征标从环面扭一下，使 Satake 参数与 Langlands 参数对齐。至于为什么只需这些迹：$\Tr(\bigwedge^d\gamma_x)$，$d=0,\ldots,n$，给出 $\gamma_x$ 的 $n$ 个特征值的全部初等对称函数，因而这 $n+1$ 个等式唯一确定半单元素 $\gamma_x$。
[来源：nLab "Tate twist"，https://ncatlab.org/nlab/show/Tate+twist（Last revised 2025-08-10；2026-11-17 访问）；并见 [Laf18] 对归一化 Satake 参数的约定（原书 p.20）。]
\end{transnote}

对幂幺根基为 $U$ 的抛物子群 $P\subsetneq\GL_n$，若球面自守函数满足
\[\int_{U(F)\backslash U(\A)}f(ug)\,du=0,\qquad g\in\GL_n(\A),\]
则称它为尖点函数。
采用通常的有限阶行列式归一化，Lafforgue 定理将具有上述非分歧特征值的尖点自守表示赋给这样的不可约 $\sigma$\cite{Laf18}。球面向量在相差一个标量的意义下唯一；Whittaker 归一化确定这个标量。
\section{Whittaker 特征标与 \texorpdfstring{$\rho$}{ρ} 移位}\label{sec:3.3}\index{Whittaker 特征标|hyperpage}\index{rho移位@$\rho$ 移位|hyperpage}
令 $N\subset\GL_n$ 为上三角幂幺子群。固定非平凡加性特征标 $\psi:k\to\Qell^\times$。在 $x$ 处，标准非退化 Whittaker 特征标通过简单根坐标得到：
\[
\begin{aligned}
\dot N(K_x)&\longrightarrow(\dot N/[\dot N,\dot N])(K_x)
\xrightarrow{\sim}\bigoplus_{\alpha\in\Delta}K_x\otimes\Omega_x\\
&\xrightarrow{\mathrm{Res}}\bigoplus_{\alpha\in\Delta}k(x)
\xrightarrow{\Tr_{k(x)/k}\circ\Sigma}k\xrightarrow{\psi}\Qell^\times.
\end{aligned}
\]
这里 $\Omega_x$ 是 $x$ 处完备化的典范线。所得特征标记为
\[\Psi_x:\dot N(K_x)\longrightarrow\Qell^\times,\]
其中上点表示 $\rho$ 移位。对 $G=\GL_n$，局部可以写为
\[\dot G_n(K_x)=\{(a_{ij})\mid a_{ij}\in K_x\otimes\Omega_x^{\otimes(j-i)}\}.\]
\srcpage{21}
因此，第一条上副对角线上的元素属于 $K_x\Omega_x$，可以对其简单根坐标取留数。$\dot G(K_x)$ 上的 Whittaker 函数是满足
\[W(ug)=\Psi_x(u)W(g),\qquad u\in\dot N(K_x)\]
的函数 $W$。通常还要求右 $\dot G(\cO_x)$ 不变性。
\section{局部 Whittaker 特征函数}\label{sec:3.4}\index{Casselman–Shalika 公式|hyperpage}
设 $\gamma\in\GL_n(\mathbb Q_\ell)$，例如上述归一化 Satake 参数 $\gamma_x$。在 $\dot G(K_x)/\dot G(\cO_x)$ 上存在唯一归一化的 Hecke 特征 Whittaker 函数 $W_{\gamma,x}$，满足：
\begin{enumerate}
\item $W_{\gamma,x}(1)=1$，其中 $1$ 表示仿射 Grassmann 簇的基点；
\item 对 $h\in\dot G(\cO_x)$，$W_{\gamma,x}(gh)=W_{\gamma,x}(g)$；
\item 对 $u\in\dot N(K_x)$，$W_{\gamma,x}(ug)=\Psi_x(u)W_{\gamma,x}(g)$；
\item 对每个 $0\le d\le n$，$T_x^dW_{\gamma,x}=q_x^{-d(d-1)/2}\Tr(\bigwedge^d\gamma)W_{\gamma,x}$。
\end{enumerate}
选取用来书写 $t_x^\lambda$ 的局部坐标后，Whittaker 双商以余权格为指标：
\[\dot N(K_x)\backslash\dot G(K_x)/\dot G(\cO_x)\simeq\Lambda,\qquad\dot N(K_x)t_x^\lambda\dot G(\cO_x)\longleftrightarrow\lambda.\]
Casselman–Shalika 公式明确给出这些值；其几何形式是 Satake 的 Whittaker 实现\cite{FGKV98}。对余权 $\lambda=(\lambda_1,\ldots,\lambda_n)$，
\[W_{\gamma,x}(t_x^\lambda)=\begin{cases}q_x^{-\sum_i(i-1)\lambda_i}\Tr(\gamma,V^\lambda),&\lambda\text{ 优势},\\0,&\text{否则},\end{cases}\]
其中，对优势 $\lambda$，$V^\lambda$ 表示最高权为 $\lambda$ 的不可约 $\checkG$ 表示。在范畴层面，这是几何 Satake 的 Whittaker 形式：
\[\operatorname{Perv}^{\mathrm{Wh}}(\dot{\Gr}_{G,x})\simeq\Rep(\checkG),\]
在导出层面则为 $D^{\mathrm{Wh}}(\dot{\Gr}_{G,x})\simeq D^b(\Rep(\checkG))$。
\section{整体 Whittaker 函数与自守函数}\label{sec:3.5}
局部特征标相乘得到整体特征标
\[\Psi=\prod_x\Psi_x:\dot N(\A)\longrightarrow\Qell^\times,\]
由留数定理，它在 $\dot N(F)$ 上平凡。上述局部定理给出定义在 $(\dot N(\A),\Psi)\backslash\dot G(\A)/\dot G(\cO)$ 上的限制乘积
\[W_\sigma:=\prod_x'W_{\gamma_x,x}.\]
$\dot G(\A)$ 上存在唯一归一化的非分歧 Hecke 特征 Whittaker 函数，由 $W(1)=1$ 及以下条件刻画：
\begin{enumerate}
\item 对 $h\in\dot G(\cO)$，$W(gh)=W(g)$；
\item 对 $u\in\dot N(\A)$，$W(ug)=\Psi(u)W(g)$；
\item $T_x^dW=q_x^{d(n-d)/2}\Tr(\bigwedge^d\sigma(\Fr_x))W$。
\end{enumerate}

对 $\GL_n$，Whittaker 展开首先给出 mirabolic 商上的尖点函数。令
\[Q_1=\begin{pmatrix}\GL_{n-1}&V_{n-1}\\0&1\end{pmatrix}\]
为 mirabolic 子群。典范映射为
\srcpage{22}
\[
\begin{aligned}
\varphi:C^\infty((\dot N(\A),\Psi)\backslash\dot G(\A))&\xrightarrow{\sim}C^\infty(\dot Q_1(F)\backslash\dot G(\A))_{\mathrm{cusp}},\\
W&\longmapsto\left(g\longmapsto\sum_{\gamma\in\dot N(F)\backslash\dot Q_1(F)}W(\gamma g)\right).
\end{aligned}
\]
等价地，该和就是相应离散商上的积分。对归一化的 Hecke 特征函数 $W_\sigma$，Piatetski-Shapiro–Shalika 自守性定理说明，$\varphi(W_\sigma)$ 实际上左 $\dot G(F)$ 不变；这就是尖点 Hecke 特征函数 $f_\sigma$。
\section{几何预览：Whittaker 范畴}\label{sec:3.6}\index{Whittaker 范畴|hyperpage}\index{手征积分|hyperpage}
上述局部陈述具有以下几何形式\cite{Gai18}：
\[\operatorname{Wh}(\dot{\Gr}_{G,x})\simeq\Rep(\checkG).\]
在曲线上取手征积分后，预期有
\[\int_X^{\mathrm{ch}}\operatorname{Wh}(\dot{\Gr}_{G,x})\simeq\int_X^{\mathrm{ch}}\Rep(\checkG).\]
左边是对阿代尔 Whittaker 商的层论替代，右边则映到 $\checkG$-局部系统叠上的拟凝聚层。图景为
\origdiagram{p022-diagram.pdf}{0.6648}
因此，Whittaker 模型是在 $\Bun_G$ 上的 Hecke 特征层与 $\LocSys_{\checkG}$ 上的层之间建立联系的局部到整体桥梁。

为进行整体比较，以 $K=k(X)$ 表示整体函数域，$\A$ 表示其阿代尔环。在函数论陈述中 $k$ 有限；几何上，将 $\psi$ 替换为 Artin–Schreier 层。令 $G=\GL_n$，$N\subset B\subset G$ 为标准的极大幂幺子群和 Borel 子群，并固定非平凡加性特征标 $\psi:k\to\Qell^\times$。为简化记号，此处省去上点：只要出现 Whittaker 特征标 $\Psi$，阿代尔群就是上文所述的 $\rho$ 扭曲形式。令
\[Q_1=\begin{pmatrix}\GL_{n-1}&V\\0&1\end{pmatrix}\]
为 mirabolic 子群。经典的目标是 Whittaker 展开映射，这一目标在文献\cite{FGKV98,Gai18}的 Whittaker 形式体系中被几何化：
\[\Av_!:C^\infty((N(\A),\Psi)\backslash G(\A))\xrightarrow{\sim}C^\infty(Q_1(K)\backslash G(\A))_{\mathrm{cusp}}.\]
它将整体 Whittaker 函数 $W_\sigma=\prod_x'W_{\gamma_x,x}$ 映为 mirabolic 商上的尖点函数 $f_\sigma$。这里 $\gamma_x$ 是上文引入的归一化 Satake 参数，局部函数 $W_{\gamma_x,x}\in C^\infty((N(K_x),\Psi_x)\backslash G(K_x)/G(\cO_x))$ 是
\srcpage{23}
Casselman–Shalika 定理给出的归一化局部 Whittaker 函数。剩下的自守性陈述是：所得 $f_\sigma$ 不仅 $Q_1(K)$ 不变，而且实际 $G(K)$ 不变。

等价地，由于 $N\subset Q_1$，有
\[(\Av_!W)(g)=\sum_{\gamma\in N(K)\backslash Q_1(K)}W(\gamma g).\]
在函数论情形，尖点性使该求和收敛。几何上，以一般位置约化模空间之间相应映射的 $!$ 推前替代它。
\section{一般位置约化与 Whittaker 特征标}\label{sec:3.7}\index{一般位置约化|hyperpage}
令 $\Bun_G^{B\text{-gen}}$ 和 $\Bun_G^{N\text{-gen}}$ 分别表示配备一般位置定义的 $B$ 约化和 $N$ 约化的 $G$ 丛预叠。有交换图
\origdiagram{p023-square.pdf}{0.3142}
这里，$\mathrm{str}$ 记住底层 $G$ 丛，$L$ 忘去一般位置的 $N$ 结构。一个 $B$ 丛的一般位置 $N$ 约化，等价于它诱导的 $T$ 丛的一般位置平凡化，因此
\[\Bun_B^{N\text{-gen}}=\Bun_B\times_{\Bun_T}\Bun_T^{\mathrm{rat\ triv}}\simeq\Bun_B\times_{\Bun_T}\operatorname{Div}_\Lambda,\]
其中 $\operatorname{Div}_\Lambda$ 是余权幺半群的带色除子空间。这是在几何上替代允许有理约化并记录其缺陷除子的做法。

\begin{transnote}{含糊点拨}
$\Bun_T^{\mathrm{rat\ triv}}$ 指带有有理平凡化的 $T$-丛：除 $T$-丛 $\mathcal L$ 外还带一个在一般点处取得的同构 $\mathcal L_\eta\simeq\mathcal L^0_\eta$。这样的对象在同构意义下由「在哪些点、以什么余权偏离平凡」完全决定，于是与 $\Lambda$-着色的除子空间 $\operatorname{Div}_\Lambda$（有限个点各标一个 $\lambda\in\Lambda$）等同；$\Lambda$ 这里取余权幺半群。所谓缺陷除子就是这组数据：它记录一般位置约化在每一点「不平凡到什么程度」，而 $\Bun_T^{\mathrm{rat\ triv}}$ 与 $\operatorname{Div}_\Lambda$ 的等同说明这部分信息可以完全交给除子承担——这也是下面 Whittaker 层能按 $\operatorname{Div}_\Lambda$ 逐点定义的原因。
[依据：原书 p.23–24 对 $\Bun_B^{N\text{-gen}}$ 与 $\operatorname{Div}_\Lambda$ 的描述；并见 [Gai18] 对一般位置约化与 Whittaker 范畴的构造。]
\end{transnote}

对 $G=\GL_n$，这里以 $\Bun_N$ 表示具有固定 $T$ 丛的 $\rho$ 扭曲 Whittaker 叠：
\[\Bun_N=\{0=\overline F_0\subset\overline F_1\subset\cdots\subset\overline F_n,\quad\overline F_i/\overline F_{i-1}\simeq\Omega_X^{n-i}\}.\]
简单根扩张类给出映射
\[\ell:\Bun_N\longrightarrow\prod_{i=1}^{n-1}\Ga\xrightarrow{\Sigma}\Ga.\]
事实上，对 $1\le i\le n-1$，第 $i$ 个相邻扩张
\[0\longrightarrow\Omega_X^{n-i}\longrightarrow\overline F_{i+1}/\overline F_{i-1}\longrightarrow\Omega_X^{n-i-1}\longrightarrow0\]
定义 $\Ext^1(\Omega_X^{n-i-1},\Omega_X^{n-i})\simeq H^1(X,\Omega_X)\simeq k$ 中的一个元素，而 $\ell$ 是这些类的和。沿 $\ell$ 拉回 Artin–Schreier 层 $\cL_\psi$，得到整体 Whittaker 层。以 $D(\Bun_G^{N\text{-gen}})^{\mathrm{Wh}}$ 表示带有这种扭曲 $N(\A)$ 等变性的范畴。

因此，对 $G=\GL_n$，mirabolic Whittaker 展开的几何形式为
\[\Av_!:D(\Bun_G^{N\text{-gen}})^{\mathrm{Wh}}\xrightarrow{\sim}D(\Bun_G^{Q_1\text{-gen}})_{\mathrm{cusp}}.\]
\srcpage{24}
它应当与局部 Whittaker–Satake 等价以及谱一侧一起理解：
\origdiagram{p024-diagram.pdf}{0.6489}
这里 $\delta_E$ 是局部系统 $E$ 处的摩天层，左边的手征积分是局部 Whittaker 范畴的整体化。
\section{退化特征标与抛物子群}\label{sec:3.8}\index{Whittaker 特征标!退化|hyperpage}
$N(\A)$ 上在 $N(K)$ 上平凡且经由简单根商分解的特征标，由简单根坐标控制：
\[N(K)^\perp:=\{\chi:N(\A)\to\Qell^\times\mid\chi|_{N(K)}=1\}_{\mathrm{sr}}\simeq K^{\oplus(n-1)}.\]
这里“sr”表示简单根商。环面 $T(K)$ 通过简单根作用，所有坐标非零的开子集等同于 $T_{\mathrm{ad}}(K)$。因此，$T(K)$ 轨道由哪些简单根坐标非零来标记。等价地，
\[\{N(K)^\perp\text{ 中的 }T(K)\text{ 轨道}\}\longleftrightarrow\{J\subset\operatorname{Dyn}(G)\}\longleftrightarrow\{P\mid B\subset P\subset G\},\]
故有 $2^{n-1}$ 个轨道。对 Levi 子群为 $M$ 的抛物子群 $P$，令 $J(P)$ 为 $M$ 的简单根集合，以 $\Psi_P:N(\A)\to\Qell^\times$ 表示一个代表特征标，其简单根坐标恰好在 $J(P)$ 上非零。一般位置特征标为 $\Psi=\Psi_G$，零特征标则是与 $P=B$ 相联系的轨道。

扩展的 Piatetski-Shapiro–Shalika 定理使用所有这些 Fourier 系数；这是文献\cite{FGKV98}中 Whittaker 比较背后的函数论输入。现在令
\[Q=\begin{pmatrix}\GL_{n-1}&V\\0&\Gm\end{pmatrix}\]
为型 $(n-1,1)$ 的极大抛物子群。存在 Fourier 系数映射
\[C^\infty(Q(K)\backslash G(\A))\longrightarrow\prod_{P\supset B}C^\infty((Z_M(K)\ltimes N(\A),\Psi_P)\backslash G(\A)),\]
其中 $Z_M$ 是 $P$ 的 Levi 子群的中心，并稳定 $\Psi_P$；该记号表示 $Z_M(K)$ 对特征标的作用平凡。限制到尖点部分后，所有退化分量都为零，并与 $P=G$ 的因子等同，即非退化 Whittaker 系数。从 $Q$ 转到 $Q_1$ 后，这就恢复了上述 mirabolic 定理。
\section{沿 mirabolic 子群迭代 Fourier 变换}\label{sec:3.9}\index{Fourier 变换!迭代|hyperpage}
在迭代证明中使用 mirabolic 子群 $Q_1=L_1\ltimes V$，其中 $L_1\simeq\GL_{n-1}$，而不使用完整的极大抛物子群。$V(K)\backslash V(\A)$ 上的 Fourier 变换使用
\[V(K)^\perp=\operatorname{Ch}(V(\A)/V(K)).\]

\begin{transnote}{易错预警}
$\operatorname{Ch}(-)$ 是特征标群（Pontryagin 对偶）的记号，全书仅此处出现：$V(\A)$ 是自对偶的，$V(K)\backslash V(\A)$ 紧且完全不连通，其对偶由 $V(\A)$ 上在 $V(K)$ 上平凡的连续特征标组成，仍以 $V(\A)/V(K)$ 为参数。因此 $V(K)^\perp$ 不是「正交补空间」，而是消没子，等于 $\operatorname{Ch}(V(\A)/V(K))$。正因为要数的是对偶这一边的轨道，才轮到 $L_1(K)$ 在 $V(K)^\perp$ 上作用——在 $V(\A)/V(K)$ 上看不出这里的分裂。
[来源：nLab "Pontrjagin dual"，https://ncatlab.org/nlab/show/Pontryagin+duality（2026-11-17 访问）；并见原书 p.24 对 $V(K)^\perp$ 的使用。]
\end{transnote}
$L_1(K)$ 在 $V(K)^\perp$ 上的作用有两个轨道：
\srcpage{25}
\begin{enumerate}
\item 经过非零特征标 $\eta_{n-1}$ 的开轨道，其稳定子是较小的 mirabolic 子群 $Q'_1(K)\subset\GL_{n-1}(K)$；
\item 闭轨道 $\{0\}$，其稳定子为 $L_1(K)$。
\end{enumerate}
因此，Fourier 变换给出分解
\[
\begin{aligned}
C^\infty(Q_1(K)\backslash G(\A))\simeq{}&C^\infty((Q'_1(K)\ltimes V(\A),\eta_{n-1})\backslash G(\A))\\
&\times C^\infty(L_1(K)\ltimes V(\A)\backslash G(\A)).
\end{aligned}
\]
若 $n\ge3$，则 $Q'_1=L'_1\ltimes V'$ 是 $\GL_{n-1}$ 中的 mirabolic 子群。第一个因子又是秩 $n-1$ 的 Whittaker 型空间，因为 $Q'_1(\A)$ 作用于 $C^\infty((V(\A),\eta_{n-1})\backslash G(\A))$。第二个因子包含在 $C^\infty(Q'_1(K)\ltimes V(\A)\backslash G(\A))$ 中，并被尖点性消去。迭代这种降低秩的 Fourier 变换后，只剩下一般位置 Whittaker 系数。

\begin{transnote}{含糊点拨}
$Q'_1\subset\GL_{n-1}$ 是 $\GL_{n-1}$ 的 mirabolic 子群，即稳定某个非零线性泛函的子群，由这里的 $\eta_{n-1}$ 给出；它此前未出现，正是下面所说的 $Q'_1=L'_1\ltimes V'$。被跳过的一步是 $\GL_{n-1}(K)$ 在 $V(K)^\perp$ 上的轨道分类：非零元素构成一条开轨道，稳定子就是稳定该泛函的 mirabolic 子群；$0$ 单独是闭轨道，稳定子为整个 $\GL_{n-1}(K)=L_1(K)$。Fourier 变换按这两条轨道把函数分成两块——开轨道给出秩降一的 Whittaker 型空间，闭轨道给出在 $L_1(K)$ 上几乎不变的部分，被尖点性消去。迭代这一步就是逐次降秩的过程。
[依据：原书 p.25 对两条轨道与 $Q'_1$ 的说明；并见 [FGKV98]（JAMS 11 (1998), 451–484）中 mirabolic 子群与 Whittaker 展开的用法。]
\end{transnote}

对于极大抛物子群 $Q$，其范畴类比表现为到粘合的扩展 Whittaker 范畴的全忠实函子：
\[D(\Bun_G^{Q\text{-gen}})\hookrightarrow\sum_{P\supset B}D(\Bun_G^{Z_M\text{-gen},\mathrm{Wh}_P}).\]
由 $P$ 标记的部分，将具有特征标 $\Psi_P$ 及 $Z_M(K)$ 约化的退化系数几何化。对于 $G=\GL_n$ 以及型为 $(n-1,1)$ 的 $P$，这是经典迭代 Fourier 变换证明的几何体现。
\begingroup\setlength{\belowdisplayshortskip}{2pt}
\section{使用凝聚层的紧化变体}\label{sec:3.10}\index{凝聚层!紧化|hyperpage}
采用使 $Q_1$ 稳定所选向量的右作用约定，$G=\GL_n$ 的 $Q_1$-一般位置约化可以写成一个非零有理向量：
\[\Bun_G^{Q_1\text{-gen}}=\{(M,\Omega_X^{n-1}\dashrightarrow M),\quad\text{不恒为 }0\}.\]
其正则版本为
\[\Bun'_n=\{(M,s:\Omega_X^{n-1}\to M)\mid s\text{ 在一般点非零},\ M\in\Bun_n\},\]
并有遗忘映射 $\Bun'_n\to\Bun_n$。由于 $M$ 局部自由且 $X$ 为光滑曲线，在一般点非零的映射 $s$ 作为凝聚层的映射是单射，并给出
\[0\longrightarrow\Omega_X^{n-1}\longrightarrow M\longrightarrow M'\longrightarrow0.\]
这里 $M'$ 是秩 $n-1$ 的凝聚层；由于 $s$ 的像未必饱和，它未必局部自由。扩张类属于 $\Ext^1(M',\Omega_X^{n-1})$，Serre 对偶给出
\[\Ext^1(M',\Omega_X^{n-1})^*\simeq\Hom(\Omega_X^{n-2},M').\]
这恰好是下一次 Fourier 变换所在的对偶向量空间。

引入紧化叠
\[\Coh'_m=\{(\cF,\Omega_X^{m-1}\xrightarrow{s}\cF)\mid\cF\in\Coh,\ \rank\cF=m\},\qquad{}^{\ne0}\Coh'_m\subset\Coh'_m,\]
其中开子叠由 $s$ 在一般点非零定义。更精确地，对 $m=n,n-1,\ldots,2$，令 $\mathcal E_m\to\Coh_{m-1}$ 为相对向量叠，它在 $M'$ 上的纤维为 $\Ext^1(M',\Omega_X^{m-1})$。Serre 对偶将对偶相对向量叠等同于 $\mathcal E_m^\vee\simeq\Coh'_{m-1}$，其在 $M'$ 上的纤维为 $\Hom(\Omega_X^{m-2},M')$。映射 $\mathcal E_m\to\Coh'_m$ 记住带有典范截面的中间项；置 ${}^{\ne0}\mathcal E_m:=\mathcal E_m\times_{\Coh'_m}{}^{\ne0}\Coh'_m$。相继的 Fourier 变换可概括为
\[{}^{\ne0}\Coh'_m\xleftarrow{p_m}{}^{\ne0}\mathcal E_m\xrightarrow{\mathsf F}\mathcal E_m^\vee\simeq\Coh'_{m-1},\qquad m=n,n-1,\ldots,2.\]
\srcpage{26}
\enlargethispage{\baselineskip}
第一次变换按照文献\cite{KL85}的归一化，定义 Fourier–Deligne 函子
\[\mathsf F_n:D({}^{\ne0}\Coh'_n)\longrightarrow D(\Coh'_{n-1}).\]
重复这个构造，每次将秩降低一步。因此，从 ${}^{\ne0}\Coh'_n$ 上的 $E$-Hecke 特征层出发，便得到 ${}^{\ne0}\Coh'_1$ 上的层。秩一的终点必定是 Whittaker 层；这是 Laumon 构造中所用 Whittaker 归一化的紧化形式。
\endgroup
